% Letter of Expectations (LOE) Template — Warner School Field Experience / Student Teaching
% Designed to match the handbook-required sections (I–VII) + Cover Letter form.

\documentclass[10pt]{article}

\usepackage[margin=0.7in]{geometry}
\usepackage{hyperref}
\usepackage{longtable}
\usepackage{array}
\usepackage{enumitem}
\usepackage{parskip}
\usepackage{titlesec}
\usepackage{lmodern}

\hypersetup{
  colorlinks=true,
  linkcolor=blue,
  urlcolor=blue,
  citecolor=blue
}

\titleformat{\section}{\large\bfseries}{\thesection}{0.75em}{}
\titleformat{\subsection}{\normalsize\bfseries}{\thesubsection}{0.75em}{}

% --- Participant information ---
\newcommand{\CandidateName}{\textbf{Piter Garcia}}
\newcommand{\CandidatePhone}{646-628-8300}
\newcommand{\CandidateEmail}{pgarcia8@u.rochester.edu}

\newcommand{\SchoolBasedTeacherName}{Derek Romig}
\newcommand{\SchoolBasedTeacherPhone}{TBD}
\newcommand{\SchoolBasedTeacherEmail}{derek.romig@greececsd.org}

\newcommand{\InclusionSBTEName}{N/A — inclusion supervision provided by university}
\newcommand{\InclusionSBTEPhone}{N/A}
\newcommand{\InclusionSBTEEmail}{N/A}

\newcommand{\SchoolNameAddress}{Pine Brook Elementary, Greece Central School District, Rochester, NY}
\newcommand{\SchoolPhone}{TBD}

\newcommand{\UnivBasedTeacherName}{Zenon (Zen) Borys — Computer Science}
\newcommand{\UnivBasedTeacherPhone}{585-259-9598}
\newcommand{\UnivBasedTeacherEmail}{zborys@warner.rochester.edu}

\newcommand{\InclusionUnivBasedTeacherName}{Sue Maddamma — Inclusion / Special Education}
\newcommand{\InclusionUnivBasedTeacherPhone}{TBD}
\newcommand{\InclusionUnivBasedTeacherEmail}{suemaddamma@gmail.com}

\newcommand{\AcademicTerm}{Spring 2026}
\newcommand{\EffectiveDate}{February 8, 2026}
\newcommand{\TargetSubmissionDate}{February 8, 2026 (two weeks after placement start)}


\begin{document}

% =========================
% COVER LETTER FORM (Handbook)
% =========================
\begin{center}
{\LARGE\bfseries COVER LETTER FOR LETTER OF EXPECTATIONS}\\[0.5em]
{\AcademicTerm}\\
\end{center}

\vspace{1.25em}
% \textbf{PLEASE FILL OUT ALL PARTS OF THIS FORM!}\\

\begin{tabular}{@{}p{0.50\textwidth}p{0.48\textwidth}@{}}
\textbf{Candidate name:} & \CandidateName \\
\textbf{Candidate phone number:} & \CandidatePhone \\
\textbf{Candidate email:} & \CandidateEmail \\
\\
\textbf{School-Based Teacher Educator name:} & \SchoolBasedTeacherName \\
\textbf{School-Based Teacher Educator phone number:} & \SchoolBasedTeacherPhone \\
\textbf{School-Based Teacher Educator email address:} & \SchoolBasedTeacherEmail \\
\\
\textbf{Inclusion School-Based Teacher Educator name:} & \InclusionSBTEName \\
\textbf{Inclusion School-Based Teacher Educator phone:} & \InclusionSBTEPhone \\
\textbf{Inclusion School-Based Teacher Educator email:} & \InclusionSBTEEmail \\
\\
\textbf{School name and address:} & \SchoolNameAddress \\
\textbf{School telephone number:} & \SchoolPhone \\
\\
\textbf{University-Based Teacher Educator name:} & \UnivBasedTeacherName \\
\textbf{University-Based Teacher Educator phone number:} & \UnivBasedTeacherPhone \\
\textbf{University-Based Teacher Educator email:} & \UnivBasedTeacherEmail \\
\\
\end{tabular}

\begin{tabular}{@{}p{0.59\textwidth}p{0.42\textwidth}@{}}
\textbf{Inclusion University-Based Teacher Educator name:} & \InclusionUnivBasedTeacherName \\
\textbf{Inclusion University-Based Teacher Educator phone number:} & \InclusionUnivBasedTeacherPhone \\
\textbf{Inclusion University-Based Teacher Educator email:} & \InclusionUnivBasedTeacherEmail \\
\end{tabular}
\vspace{1em}

\textbf{Acknowledgement}\\
All parties have read and agree with the objectives and requirements outlined in this Letter of Expectations and have read the entire Student Teaching Handbook.

\vspace{1.25em}

\begin{tabular}{@{}p{0.88\textwidth}p{0.18\textwidth}@{}}
\textbf{Candidate:} \rule{0.78\textwidth}{0.4pt} & \textbf{Date:} \rule{0.10\textwidth}{0.4pt} \\
\textbf{School-Based Teacher Educator:} \rule{0.58\textwidth}{0.4pt} & \textbf{Date:} \rule{0.10\textwidth}{0.4pt} \\
\textbf{Inclusion School-Based Teacher Educator:} \rule{0.49\textwidth}{0.4pt} & \textbf{Date:} \rule{0.10\textwidth}{0.4pt} \\
\textbf{University-Based Teacher Educator:} \rule{0.54\textwidth}{0.4pt} & \textbf{Date:} \rule{0.10\textwidth}{0.4pt} \\
\textbf{Inclusion Univ-Based Teacher Educator:} \rule{0.50\textwidth}{0.4pt} & \textbf{Date:} \rule{0.10\textwidth}{0.4pt} \\
\end{tabular}

\newpage

% =========================
% LETTER OF EXPECTATIONS (Handbook sections I–VII)
% =========================
\begin{center}
{\LARGE\bfseries LETTER OF EXPECTATIONS}\\[0.25em]
{\AcademicTerm}\\
\end{center}

\textbf{Effective date:} \EffectiveDate\\
\textbf{Target submission date:} \TargetSubmissionDate\\

\textbf{Purpose:} This document clarifies schedule commitments, roles, communication norms, lesson plan expectations, meeting and observation routines, and contingency procedures for the field experience and student teaching placement.

\section{General Expectations}
% \textit{What do the School-Based Teacher Educator(s) and the candidate expect to occur during the field experience and/or student teaching experience?}
This placement follows a graduated release-of-responsibility model across two phases:

\textbf{Phase 1 --- Observation \& Assistantship (150+ hours, Feb--mid-Apr 2026):}
\begin{itemize}[leftmargin=*,itemsep=0.25em]
  \item Learn the classroom context, students, routines, and IGNITE curriculum (CS, design thinking, robotics, coding).
  \item Observe instruction and document professional learning through structured notes and reflections.
  \item Support students one-on-one and in small groups; assist with lesson preparation, materials, and formative checks.
  \item Co-teach lesson segments and mini-lessons with increasing scope.
  \item Develop practice skills—differentiation, UDL-informed supports, and culturally and disability-responsive pedagogy.
\end{itemize}

\textbf{Phase 2 --- Student Teaching (5 weeks / 25 instructional days, mid-Apr--Jun 2026):}
\begin{itemize}[leftmargin=*,itemsep=0.25em]
  \item Assume primary or co-teaching responsibility for lessons and units.
  \item Plan and deliver full lesson sequences, including differentiation and inclusive supports.
  \item Implement formative and summative assessments; analyze results to adjust instruction.
  \item Manage classroom routines with support from the cooperating teacher.
  \item Collect evidence and artifacts for CPAST evaluation and reflective practice.
\end{itemize}

\textbf{Candidate learning goals:}
\begin{enumerate}[leftmargin=*,itemsep=0.2em]
  \item \textbf{Planning \& alignment:} Produce weekly plans with clear learning targets aligned to standards/curriculum.
  \item \textbf{Inclusive instruction:} Implement uniform differentiation supports (scaffolds, choice, multimodal representations).
  \item \textbf{Assessment \& adjustment:} Use at least one formative check per lesson block; document one instructional adjustment per week based on evidence.
  \item \textbf{Digital tools in CS learning:} Integrate age-appropriate tools to support coding/CS learning and collaboration.
  \item \textbf{Reflective practice:} Weekly reflection notes tying instruction to research and identifying one improvement action.
\end{enumerate}

\section{Specific Requirements}
% \textit{In as detailed language as possible, explain the duties of each participant in the field experience and/or student teaching experience.}

\subsection{Candidate}
% \textit{What will the candidate do in terms of observation, lesson planning, teaching, evaluation and assessment, parent contact, administrative contact, working with other teachers (or candidates), one-on-one student instruction, administrative duties, after-school duties, etc.?}

\begin{itemize}[leftmargin=*,itemsep=0.25em]
  \item \textbf{Observation:} Maintain structured observation notes focused on instructional strategies, classroom management, student engagement, differentiation, and the IGNITE curriculum. Document reflections weekly.
  \item \textbf{Lesson planning:} Submit a weekly planning packet to Derek by Friday for the following week. Packet includes: learning targets/objectives, lesson outline, inclusion/differentiation plan (UDL supports, scaffolds, accommodations), formative checks, materials list, and an assessment artifact (exit ticket, rubric, quiz, or reflection prompt).
  \item \textbf{Teaching:} Phase~1---co-teach lesson segments and lead small-group instruction. Phase~2---plan and deliver full lessons/units with increasing independence.
  \item \textbf{Assessment:} Use at least one routine formative check per lesson block; document at least one instructional adjustment per week based on evidence. Implement summative assessments as appropriate during Phase~2.
  \item \textbf{Parent/admin contact:} Follow school policy; all parent communication goes through Derek. Maintain strict confidentiality of student information.
  \item \textbf{Collaboration:} Work with grade-level teams, specialists, and support staff as directed by Derek. Participate in staff meetings and professional development when present.
  \item \textbf{Administrative/after-school duties:} Arrive on time for committed blocks; assist with arrival/dismissal routines as requested. No after-school duties are required, but candidate may voluntarily participate.
  \item \textbf{One-on-one instruction:} Tutor individual students or provide targeted support as identified by Derek.
  \item \textbf{Professional conduct:} Maintain professional conduct, confidentiality, boundaries, and timely communication.
  \item \textbf{Evidence portfolio:} Maintain an organized collection of lesson plans, observation notes, reflections, assessment artifacts, and de-identified student work throughout the semester.
  \item \textbf{Midpoint self-evaluation:} Complete the Warner\_MidpointEval form as a self-evaluation at approximately 75 placement hours (midpoint of Phase~1). Sit down with Derek to compare evaluations side by side and discuss perceptions of progress. The five weeks of student teaching cannot begin until this form is submitted.
  \item \textbf{Final self-assessment (CPAST):} The CPAST rubric is the candidate's final student teaching assessment. At the end of the placement, the candidate will self-assess using the CPAST rubric, then participate in a three-way conference (with Derek and Zenon) where all three parties individually rate performance, discuss ratings, work to reach consensus scores, and articulate growth goals. To support consistent scoring, all three participants complete the brief CPAST online training prior to the final conference.
\end{itemize}

\subsection{School-Based Teacher Educator (Derek Romig)}
% \textit{What will the School-Based Teacher Educator do in terms of observation, assisting with unit and lesson planning, modeling teaching, critiquing and advising, and acting as a liaison between the candidate and other teachers, administrators, and parents, etc.?}

\begin{itemize}[leftmargin=*,itemsep=0.25em]
  \item \textbf{Orientation:} Introduce Piter to classroom routines, school procedures, the IGNITE curriculum (CS, design thinking, robotics, coding), and district/building policies.
  \item \textbf{Modeling \& coaching:} Demonstrate lessons and instructional strategies; co-plan lessons with Piter; provide regular verbal and written coaching feedback.
  \item \textbf{Observations:} Observe the candidate informally weekly and formally before each university supervisor visit. Share observation notes and actionable feedback promptly.
  \item \textbf{Liaison:} Facilitate introductions to other teachers, administrators, specialists, and parents. All parent/guardian communication flows through Derek.
  \item \textbf{Midpoint evaluation:} Complete the Warner\_MidpointEval form for the candidate at approximately 75 placement hours. Sit down with Piter to compare evaluations (Derek's and Piter's self-evaluation) side by side and discuss progress. Submit Derek's copy to Zenon by email (\href{mailto:zborys@warner.rochester.edu}{\nolinkurl{zborys@warner.rochester.edu}}). Suggested filename: \nolinkurl{Garcia_Romig_Spring2026MIDPOINT.pdf}.
  \item \textbf{CPAST training \& final assessment:} Complete the CPAST online training (60--90 min, 85\% passing score required). At the end of the placement, individually rate Piter using the CPAST rubric, participate in the three-way conference to reach consensus scores, and submit Derek's completed CPAST rubric (with any notes) to Zenon.
  \item \textbf{Gradual release:} Provide increasing opportunity for the candidate to assume instructional responsibility, from co-teaching segments to full lesson leadership.
\end{itemize}

\subsection{University-Based Teacher Educators (Zenon Borys \& Sue Maddamma)}
% \textit{Explain when and how often the University-Based Teacher Educators will meet with the School-Based Teacher Educator and candidate to provide professional support, observe the candidate, and discuss the candidate's progress.}

\begin{itemize}[leftmargin=*,itemsep=0.25em]
  \item \textbf{Observation/visit cadence:} Zenon Borys (CS) will conduct 2--3 visits during the semester: one during Phase~1 (field experience) and 1--2 during Phase~2 (student teaching). Sue Maddamma (Inclusion) will coordinate with Zenon and join periodically for inclusion-focused feedback.
  \item \textbf{Feedback \& coaching:} After each observation, the supervisor(s) will meet with the candidate and Derek as soon as possible to debrief. Written observation notes and actionable feedback will be provided.
  \item \textbf{Midpoint conference:} At approximately 75 placement hours, Zenon will meet with Derek and Piter to review the midpoint evaluation (Derek's assessment + Piter's self-evaluation) and discuss progress. The five-week student teaching phase cannot begin until this form is submitted.
  \item \textbf{CPAST training \& final assessment:} The CPAST rubric is Piter's final student teaching assessment. Prior to the final conference, Zenon completes the CPAST online training (and confirms completion for all participants who will rate using the rubric). At the end of the placement, Zenon will facilitate a three-way conference (via Zoom) where Derek, Piter, and Zenon each individually rate Piter's performance, discuss ratings, work to reach consensus scores for each category, and articulate goals for Piter's growth. After the conference, Zenon submits the consensus scores as required and Derek submits his completed CPAST rubric with notes to Zenon.
  \item \textbf{Availability between visits:} Zenon is available by email (\texttt{zborys@warner.rochester.edu}) and phone (585-259-9598). Sue coordinates with Zenon. If questions or concerns emerge outside scheduled visits, either party should contact the supervisor promptly.
  \item \textbf{Documentation:} Supervisors coordinate all required documentation submission processes, including midpoint and final evaluation forms.
\end{itemize}

\section{Schedule}

\subsection{Daily Schedule Expectations}
% \textit{What time is the candidate expected to arrive? Until what time is the candidate expected to stay?}

The candidate is simultaneously enrolled in required graduate coursework at RIT and UofR. The schedule below reflects a stable weekly commitment that balances consistent placement presence, safety, and sustainable commuting, and coursework responsibilities.

\begin{center}
\begin{tabular}{|l|l|l|}
\hline
\textbf{Day} & \textbf{Placement Presence} & \textbf{Notes} \\
\hline
Monday    & 1:10 PM -- 3:10 PM  & Afternoon block (after RIT lectures) \\
Tuesday   & 9:00 AM -- 3:20 PM  & Full day \\
Wednesday & 1:10 PM -- 3:10 PM  & Afternoon block (after RIT lectures) \\
Thursday  & 9:00 AM -- 3:20 PM  & Full day \\
Friday    & 9:10 AM -- 3:10 PM  & Full day \\
\hline
\end{tabular}
\end{center}

\textbf{Weekly total (target):} $\sim$22 hours/week of placement presence.

\textbf{Check-in procedure:} Sign in at the front office upon arrival; report directly to Derek's classroom.

\textbf{Scheduling constraints:} If an accommodation is required for a specific week (e.g., special seminar dates), the candidate will communicate and confirm the plan in advance with Derek and Zenon.

\subsection{Student Teaching Commitments}
% \textit{Address the following in specific terms.}

\begin{enumerate}[leftmargin=*,itemsep=0.25em]
  \item \textbf{When will the candidate begin to assume partial or full responsibility?}\\
        Partial responsibility (co-teaching segments, small-group instruction) begins during Phase~1 and increases gradually. Full responsibility begins at the start of Phase~2 ($\sim$mid-April 2026), contingent on a passing midpoint evaluation and submission of this Letter of Expectations.
  \item \textbf{Which classes will the candidate co-teach with the SBTE?}\\
        The candidate will co-teach IGNITE program sessions (computer science, design thinking, robotics, and coding lessons) across grade levels served by Derek's classroom.
  \item \textbf{Which classes will the candidate observe?}\\
        During Phase~1, the candidate will observe all of Derek's instructional periods, including transitions, morning routines (Tue/Thu/Fri full days), and afternoon blocks (Mon/Wed). Observations focus on pedagogy, classroom management, differentiation, and student engagement.
  \item \textbf{Transition from less to more instructional responsibility:}\\
        \textbf{Weeks 1--3:} Observation, structured notes, 1:1 and small-group support.\\
        \textbf{Weeks 4--6:} Co-teach lesson segments; lead mini-lessons under Derek's guidance.\\
        \textbf{Weeks 7--10:} Increase co-teaching scope; begin planning full lessons.\\
        \textbf{Phase~2 (Weeks 11--16):} Candidate leads full lessons/units with Derek providing feedback. Candidate teaches for 25 instructional days, with options including: leading one class while Derek leads another, co-teaching together, or teaching specific lesson units independently.
\end{enumerate}

\subsection{Timeline of responsibilities}
\textit{Include a timeline of responsibilities over the course of the placement.}

\begin{longtable}{|p{0.22\textwidth}|p{0.36\textwidth}|p{0.34\textwidth}|}
\hline
\textbf{Weeks / Dates} & \textbf{Candidate responsibilities} & \textbf{Evidence / checkpoints} \\
\hline
\endfirsthead
\hline
\textbf{Weeks / Dates} & \textbf{Candidate responsibilities} & \textbf{Evidence / checkpoints} \\
\hline
\endhead
\hline
\endfoot
Feb 4--14 (Wk 1--2) & Observe instruction; take structured notes; support students 1:1 and in small groups; learn routines and IGNITE curriculum & Observation notes; initial reflections \\
\hline
Feb 17--Mar 7 (Wk 3--5) & Assist with lesson prep and materials; co-teach lesson segments; begin formative checks & 1--2 lesson plan artifacts; co-teaching log \\
\hline
Mar 10--28 (Wk 6--8) & Lead mini-lessons; increase co-teaching scope; plan full lessons with Derek's guidance & Lesson plans with differentiation; formative assessment examples \\
\hline
$\sim$Mar 15 & \textbf{Midpoint evaluation} ($\sim$75 hrs): candidate self-assessment + Derek assessment + supervisor conference & \nolinkurl{Garcia_Romig_Spring2026MIDPOINT.pdf} submitted to Zenon \\
\hline
Mar 31--Apr 10 (Wk 9--10) & Continue co-teaching with increasing independence; finalize Phase~2 transition plan & Observation feedback; revised learning goals \\
\hline
$\sim$Apr 15--Jun (Phase 2) & Candidate leads full lessons/units (25 instructional days); manages routines; implements assessments; collects CPAST evidence & Weekly lesson plans; formative/summative artifacts; portfolio evidence \\
\hline
Late Apr / May & \textbf{Final three-way conference} (Derek, Piter, Zenon): consensus CPAST ratings and growth planning & Final CPAST rubrics submitted to Zenon \\
\hline
\end{longtable}

\section{Lesson Plans}
% \textit{How often and when will the candidate review their lesson plans with the School-Based Teacher Educator(s), and what will happen during those reviews?}

\begin{itemize}[leftmargin=*,itemsep=0.25em]
  \item \textbf{Review frequency:} Weekly. The candidate submits a planning packet to Derek by \textbf{Friday end-of-day} for the following week. If a school holiday or major assignment disrupts Friday submission, Derek and the candidate agree on an alternative deadline in advance.
  \item \textbf{Review method:} Shared document (Google Doc or email attachment). Derek reviews the packet before the start of the following week and provides written or verbal feedback.
  \item \textbf{Review focus:} Alignment of learning targets to standards/curriculum; quality of differentiation and inclusive supports (UDL); clarity of formative checks for understanding; materials and technology readiness.
  \item \textbf{Planning packet minimum contents:}
  \begin{itemize}[leftmargin=1.5em,itemsep=0.15em]
    \item Learning targets/objectives (student-friendly wording where appropriate)
    \item Outline of lesson flow (opener $\to$ modeling $\to$ guided practice $\to$ independent practice $\to$ closure)
    \item Inclusion/differentiation plan (scaffolds, accommodations/modifications)
    \item Formative checks (what, when, and adjustment plan if students struggle)
    \item Materials/tech list (links, handouts, device needs)
    \item Assessment artifact (exit ticket, rubric, quiz, reflection prompt, or performance task)
  \end{itemize}
\end{itemize}

\newpage

\section{Meetings}

\subsection{Regular Meetings}
\textit{How often will the candidate and School-Based Teacher Educator(s) meet? When will these meetings take place? What will be the focus?}

\begin{itemize}[leftmargin=*,itemsep=0.25em]
  \item \textbf{Weekly check-in (15--30 minutes):} Derek and Piter meet once per week (day/time to be agreed, recommended: end of a full teaching day) to:
  \begin{itemize}[leftmargin=1.5em,itemsep=0.1em]
    \item Preview upcoming lessons and confirm roles for co-teaching/support.
    \item Address schedule adjustments for the following week.
    \item Review feedback from the prior week's instruction.
  \end{itemize}
\end{itemize}

\subsection{Additional Meetings}
\textit{Will the candidate and School-Based Teacher Educator(s) be able to meet at other times if the need arises? How will these meetings be arranged?}

\begin{itemize}[leftmargin=*,itemsep=0.25em]
  \item Yes. Either party may request an ad hoc meeting via email or text. Meetings will be arranged at a mutually convenient time, typically before or after the placement block.
  \item Day-to-day classroom questions and logistics should be directed to Derek first.
  \item For scheduling or evaluation process questions, contact Zenon Borys (CS) or Sue Maddamma (Inclusion).
\end{itemize}

\subsection{Post-Observation Meetings}
\textit{The candidate, School-Based Teacher Educator(s), and University-Based Teacher Educator will meet as soon as possible after each observation by the University-Based Teacher Educator.}

\begin{itemize}[leftmargin=*,itemsep=0.25em]
  \item After each university supervisor observation, Zenon (and/or Sue) will debrief with Derek and Piter as soon as possible---ideally the same day or within 48 hours.
  \item The debrief will review the lesson observed, discuss strengths and areas for growth, and set specific goals for subsequent instruction.
  \item Monthly: Zenon visits for observation/feedback; Sue joins periodically for inclusion-focused feedback. All parties review progress toward milestones.
\end{itemize}

\section{Observations}

\subsection{School-Based Teacher Educator Observations}
\textit{When will the School-Based Teacher Educator(s) observe the candidate (e.g., before the University-Based Teacher Educator/faculty come to observe)?}

\begin{itemize}[leftmargin=*,itemsep=0.25em]
  \item Derek will observe the candidate informally on a \textbf{weekly} basis during co-teaching and independent lesson segments.
  \item Derek will conduct a more structured, formal observation \textbf{before} each scheduled university supervisor visit so that feedback can inform the observed lesson.
  \item Feedback from SBTE observations will be shared verbally (same day) and/or in brief written notes.
\end{itemize}

\newpage
\subsection{University-Based Teacher Educator Observations}
\textit{When will the University-Based Teacher Educator observe the candidate teaching lessons during the field experience or student teaching?}

\begin{itemize}[leftmargin=*,itemsep=0.25em]
  \item \textbf{Zenon Borys (CS):} 2--3 visits total --- first observation $\sim$Feb 16--20; one observation during Phase~1; 1--2 observations during Phase~2 (student teaching).
  \item \textbf{Sue Maddamma (Inclusion):} Coordinates with Zenon; joins periodically for observation and feedback.
  \item \textbf{Evidence captured during observations:} Lesson plan for the observed session, supervisor observation notes, candidate post-lesson reflection, and de-identified student work samples where relevant.
  \item \textbf{Suggested evidence cadence (candidate maintains throughout):}
  \begin{itemize}[leftmargin=1.5em,itemsep=0.1em]
    \item 1--2 lesson plans/week saved as PDF
    \item 1 formative assessment example/week + resulting adjustment note
    \item 1 weekly reflection (what worked, what didn't, what changes next)
    \item Periodic photos of non-identifying classroom materials or student-facing instructions
  \end{itemize}
\end{itemize}

\section{Possible Situations / Contingencies}

\begin{enumerate}[leftmargin=*,itemsep=0.25em]
  \item \textbf{What will the candidate do if the SBTE is absent?}\\
        The candidate is \textbf{not legally permitted} to substitute for the SBTE. If Derek is absent, Piter will report to the assigned substitute teacher or building administrator for direction. The candidate may observe or assist the substitute but will not assume sole classroom responsibility. Piter will notify Zenon of the situation.

  \item \textbf{What will the SBTE do if the candidate is absent?}\\
        Derek will proceed with instruction as normal. Any lesson components the candidate was expected to lead will be covered by Derek or reassigned. The candidate will provide an updated plan for missed responsibilities upon return.

  \item \textbf{How will each party inform the other of absences?}\\
        The absent party will notify the other as early as possible (ideally the evening before or early morning of the absence) via text/phone for urgent same-day notice and a follow-up email for documentation. The candidate will also notify Zenon Borys of any absence.

  \item \textbf{What will be done in case of snow days or emergencies?}\\
        Follow Greece Central School District policy for closures/delays. If the district closes, no placement attendance is required. The candidate will complete an agreed-upon asynchronous task (e.g., reflection writing, lesson planning, or resource review) as arranged with Derek. Any full or partial absences that reduce the candidate's total required hours must be made up with additional time.
\end{enumerate}

\textbf{Note:} Candidates are not legally allowed to substitute for School-Based Teacher Educators during clinical experiences.

% \vfill

% \textbf{Signatures}\\
% \begin{tabular}{@{}p{0.68\textwidth}p{0.28\textwidth}@{}}
% \textbf{Candidate:} \rule{0.48\textwidth}{0.4pt} & \textbf{Date:} \rule{0.18\textwidth}{0.4pt} \\
% \textbf{School-Based Teacher Educator:} \rule{0.28\textwidth}{0.4pt} & \textbf{Date:} \rule{0.18\textwidth}{0.4pt} \\
% \textbf{University-Based Teacher Educator:} \rule{0.24\textwidth}{0.4pt} & \textbf{Date:} \rule{0.18\textwidth}{0.4pt} \\
% \end{tabular}

\end{document}